% Lösungen Einheit 2 von 4 – Punktprobe und Punkte auf der Geraden (zusammenbau v0.1)
\einheitenkopf{Einheit 2 von 4 · Punktprobe und Punkte auf der Geraden}

\erg{13}{a) aus der zweiten oder der dritten – die erste Komponente des Richtungsvektors ist null \quad b) Aus der ersten Koordinate $r = 3$; zweite: $2 + 3 = 5$ ✓, dritte: $-1 + 9 = 8$ ✓ – $P$ liegt auf $g$. \quad c) Jeder Punkt von $g$ hat die zweite Koordinate $-1$ (Richtungskomponente null), $A$ hat 3. \quad d) $1 - r = -2$ liefert $r = 3$; $y_B = 3 + 6 = 9$, $z_B = 2 - 9 = -7$ \quad e) $|\vec u| = 3$, also $A + 3 \cdot \vec u = (9 | -2 | 8)$ (oder $A - 3 \cdot \vec u = (-3 | 4 | -4)$) \quad f) Ansatz $(2 | 6 | 0) = (8 | 0 | 0) + t \cdot (-8 | 8 | 0)$: erste und zweite Koordinate liefern $t = 0{,}75$, die dritte stimmt, und $0 \le 0{,}75 \le 1$ – $T$ liegt auf der Strecke. \quad g) $\overrightarrow{AB} = (2 | 1 | -2)$, $\overrightarrow{AC} = (6 | 3 | -6) = 3 \cdot \overrightarrow{AB}$ – die Verbindungsvektoren sind kollinear, $A$, $B$, $C$ liegen auf einer Geraden. \quad h) $|\overrightarrow{AE}| = 21$ LE, 84 m sind $16{,}8$ LE, also $\lambda = 16{,}8 : 21 = 0{,}8$; $x_3 = 2 + 0{,}8 \cdot 9 = 9{,}2$ LE, Höhe: $9{,}2 \cdot 5 = 46$ m \quad i) $\vec x = (5 | 0 | 0) + r \cdot (-4 | 4 | 6)$; aus der ersten Koordinate $r = 0{,}5$, zweite: $0 + 2 = 2$ ✓, dritte: $0 + 3 = 3 \ne 5$ – $S$ liegt nicht auf der Geraden. \quad j) Ansatz: alle drei Koordinaten gleich null. Aus der zweiten Koordinate $r = 2k$, in die erste eingesetzt $2 + 4k - 4k = 2 \ne 0$ – Widerspruch für jedes $k$.}
\erg{14}{a) $J = (2 + 0{,}9t | 1 + 1{,}2t | 3)$ mit $0 \le t \le 1$; $|\overrightarrow{IJ}|^2 = 2{,}25t^2 + 1{,}44 = 2{,}25$ liefert $t = 0{,}6$ ($t = -0{,}6$ liegt nicht auf dem Balken); Verhältnis $0{,}6 : 0{,}4 = 3 : 2$}
\erg{15}{a) Die dritte Koordinate fehlt: $2 + 2 = 4 \ne 7$ – Widerspruch, $P$ liegt nicht auf $g$. \quad b) Ein Punkt liegt nur dann auf der Geraden, wenn derselbe Parameterwert alle drei Koordinaten liefert; liefert eine Koordinate einen anderen Wert, gibt es keinen solchen Parameter. \quad c) Höhe $1\,200t = 900$ liefert $t = 0{,}75$; ab $(225 | 300 | 900)$ ist er höher als 900 m.}
